\section{Related Work}
\label{sec:related}

\subsection{Honeypot and Honeynet Architectures}
Deception research has moved from single-service traps to orchestrated
honeynets~\cite{spitzner2002}, with recent surveys cataloging managed
deception fabrics and cloud-native deployment~\cite{lopez2024}. Multi-sensor
platforms aggregate dozens of low- and medium-interaction sensors in one
stack, and containerized honeynets scale the same idea
elastically~\cite{t-pot2024}. The survey's gap analysis is blunt: no
widely deployed framework grounds its responses in real-time threat
intelligence or differentiates behavior across attacker classes, which
are the two capabilities G1 and \S\ref{sec:chrollo} exist to provide.
What these architectures share is a static
relationship between the decoy and its content: the deception is authored
once, at deployment time. \system differs in kind, not degree: its served
content is decided per session, and every decision passes a verification
gate before reaching the wire.

\subsection{Generative Models in Deception}
LLM-backed honeypots are the closest line of work, and the field has
moved quickly from feasibility demonstrations to state-aware designs.
GPT-driven SSH bots demonstrate that language models can hold convincing
interactive sessions~\cite{kouremetis2024}; LLM-generated web deception
serves convincing HTTP responses for arbitrary paths~\cite{galah2024};
RAG-coupled generation grounds responses in retrieved
content~\cite{qiao2025,sezgin2025decoypot}, and agentic RAG pipelines
have been proposed for CTI workflows~\cite{belmahjoub2025,gao2025}.
HoneyGPT breaks the classical flexibility/depth/deception trilemma with
structured prompting and reports a three-month field
evaluation~\cite{wang2026honeygpt}; DecoyPot drives a web-API honeypot
with a RAG agent improved from live attacker
traffic~\cite{sezgin2025decoypot}; and LLMHoney pairs a virtual file
system with real-time LLM response generation~\cite{malhotra2025llmhoney}.
Most directly relevant to this paper's machinery, a 2026 generation of
systems now treats state and validation explicitly: LLM-SSHH serializes a
persistent state snapshot into the prompt and validates generated
responses against state constraints, rejecting hallucination-driven
inconsistencies~\cite{li2026llmsshh}, and GenPot couples generative
interaction across command-line, API, and web surfaces with state
management and guardrails, evaluated operationally and at
scale~\cite{laragutierrez2026genpot}. A recent SoK systematizes this
design space and its failure modes~\cite{bridges2026}.

The gap this paper addresses is therefore no longer whether LLM honeypots
can be consistent -- the newest systems demonstrate exactly that -- but
what the consistency machinery is \emph{for}. Prior systems improve or
validate generative consistency: state snapshots and constraint checks
(LLM-SSHH), guardrails and state management (GenPot), structured
prompting (HoneyGPT) are mechanisms for making the generator's output
stay believable. This work turns the entire serving path into an
explicitly verified security architecture: every response is either
admitted by a serving-time verifier against a session-scoped fiction
graph with a recorded verdict, served from human-vetted deterministic
content, or withheld -- with the verifier itself red-teamed by labeled
cases and live-LLM-authored drafts (\S\ref{sec:rq2}) -- and every safety
decision lands in an append-only audit trail a human gate can act on
(G3, G5). The second difference is methodological and, we argue, the
paper's central contribution: the intelligence-ladder evaluation places
every mechanism on a certifiability axis and prices each rung by measured
marginal effect, including the fold-down direction -- an adaptive rung
whose benefit a certifiable mechanism reproduces is removed from the
serving path (\S\ref{sec:jeveval}). Prior systems argue \emph{for} their
intelligent components; this evaluation is designed to argue
\emph{against} them and to keep only what survives pricing.
Table~\ref{tab:related-matrix} summarizes the positioning along the
dimensions this paper treats as first-class security properties.

\begin{table*}[t]
\centering
\caption{Positioning against recent generative-honeypot systems. Y =
present as a first-class, evaluated mechanism; P = partial (mechanism
present but not evaluated under adversarial pressure, or scoped to
response content only); N = not reported. Entries reflect the cited
papers' own evaluations; where a capability is absent from a paper's
reported design we mark N without implying the system could not add it.
\system's column is the only one whose evaluation includes a priced
fold-down of an intelligent component.}
\label{tab:related-matrix}
\small
\setlength{\tabcolsep}{3.4pt}
\begin{tabular}{lccccccccc}
\toprule
 & \rotatebox{60}{State consistency} & \rotatebox{60}{Serving-time verification} & \rotatebox{60}{Danger/safety gate} & \rotatebox{60}{Deterministic fallback} & \rotatebox{60}{CTI grounding} & \rotatebox{60}{Human-in-the-loop gate} & \rotatebox{60}{Verifier red-team} & \rotatebox{60}{Field deployment} & \rotatebox{60}{Mechanism fold-down} \\
\midrule
HoneyGPT~\cite{wang2026honeygpt} & P & N & N & N & N & N & N & Y & N \\
DecoyPot~\cite{sezgin2025decoypot} & P & N & N & N & Y & N & N & N & N \\
LLMHoney~\cite{malhotra2025llmhoney} & P & N & N & P & N & N & N & N & N \\
LLM-SSHH~\cite{li2026llmsshh} & Y & P & N & N & N & N & N & N & N \\
GenPot~\cite{laragutierrez2026genpot} & P & P & P & N & N & N & N & P & N \\
\system (this work) & Y & Y & Y & Y & Y & Y & Y & P & Y \\
\bottomrule
\end{tabular}
\end{table*}

\subsection{LLM Security and Alignment Pressure}
Our provider-fragility observations connect to a broader alignment literature.
Safety-aligned models refuse or deflect on adversarial-role prompts, and
constitutional and RLHF training shapes exactly the refusal behavior our
harness must route around~\cite{bai2022,ganguli2022,christiano2017}.
Prompt-injection risks in retrieval pipelines motivate the corpus hygiene
of \S\ref{sec:implementation}~\cite{greshake2023,perez2022}, and
documented failures of ML-backed security tooling under distribution
shift~\cite{reworr2024} reinforce the case for measured fallbacks over
trusting generation. The generate-then-check pattern itself has
classical roots, and modern guardrail and safety-filtering systems
structure the same draft-then-verify separation around content policy
rather than world-model consistency~\cite{inan2023}; our verifier
differs in checking against a session-scoped fiction the defender fully
controls, which makes the check a decision procedure rather than
another learned judgment. The deployment lesson we add is operational:
treat the provider as an unreliable component with catalog drift, quota
exhaustion, and alignment-induced timeouts as normal operating conditions.
The reviewer-facing implication is a methodological one: evaluations of
LLM-driven defenses should report provider instability as part of the
results, because it shaped which experiments were even runnable.

\subsection{CTI and Behavior Analysis}
Hybrid RAG over threat intelligence has been shown effective for
attribution and technique mapping~\cite{alsada2024,rajapaksha2024},
agentic RAG for attack investigation and incident
response~\cite{blefari2025,liu2025}, and TTP extraction from unstructured
text predates the current generation of systems~\cite{orbinato2022}.
RAG-driven deception in the honeypot context specifically is represented
by DecoyPot's self-improving RAG agent~\cite{sezgin2025decoypot} and
adaptive RAG honeypot architectures~\cite{qiao2025}; what these systems
ground is response content, whereas this paper's retrieval rung grounds
the \emph{plan} (which facts to project, which bait to attach to which
predicted technique) and its effect is measured by a paired ablation with
a structural-dependence control (\S\ref{sec:rq1}).
Behavioral profiling of attacker sessions underpins both intrusion
detection~\cite{bachl2021,farasat2024b,sharafaldin2018} and
honeypot-facing classification~\cite{lekssays2025}, and attack-graph
sequence models connect observed actions to
attribution~\cite{alsaheel2021,noureddine2020}. Adaptive
deception experiments show that attacker skill modulates effective
strategy~\cite{kouremetis2024,jeon2024}, and honeytoken-based
attribution~\cite{spitzner2003,valeros2024} motivates the
canary-registered bait
design. \system composes these primitives into a loop whose safety
properties, rather than its raw capability, are the contribution.
